\scanpage{13}
\mspage{5}

qui avec l'origine $e$ forment l'ens
des pts fixes de $X_S$. (L'ens
des pts fixes sous $\mathbb Z/6\mathbb Z$ \unsure{est|se} réduit
à $e$.) Les deux pts fixes
sous $\mathbb Z/3\mathbb Z$, distincts de $e$, corres-\newline
pondent aux deux triangles coinscrits
\unsure{de|dans} l'hexagone $S$. (\inserted{les deux} triplets \unsure{des|de}
cotés mutuellement non adjacents.)

D\uline{écompositions cellulaires}\newline
ca\uline{noniques}. Décomposition en
deux faces triangulaires, par trois
arêtes, un seul sommet. Les
faces triangulaires correspondent aux
éléments de $\unsure{(X_S)^\zeta|(X_S^\zeta)}\setminus\{e\}$ (chaque
face en contient un et un seul) i.e.
aux triangles coinscrits de l'hexagone-\newline
Les arêtes correspondent aux trois

\margin{12}
